\begin{table}[htbp]
\centering
\caption{County-level border effect and entry-cap interaction (CO, OR, MT, NM; quarterly)}
\label{tab:county_border_capped_main}
\footnotesize
\begin{tabular}{lccccc}
\toprule
& \multicolumn{2}{c}{Border} & \multicolumn{2}{c}{Border $\times$ Capped} & \\
Outcome & Coef. (\%) & $p$ (cl.) & Coef. (\%) & $p$ (cl.) & $n$ \\
\midrule
Revenue per capita & $+0.860$ ($+136.4\%$) & 0.000 & $-0.720$ ($-51.3\%$) & 0.015 & 3191 \\
Revenue per store & $+0.619$ ($+85.7\%$) & 0.121 & $-0.439$ ($-35.5\%$) & 0.232 & 1720 \\
Dispensary density & $+0.241$ ($+27.3\%$) & 0.451 & $-0.126$ ($-11.9\%$) & 0.618 & 1720 \\
\bottomrule
\end{tabular}
\begin{minipage}{0.98\textwidth}
\vspace{0.3em}
\footnotesize
\textit{Notes:} Sample: CO, OR, MT, NM (Washington dropped). Border = county centroid within 75 miles of the
nonlegal-state line, with $\geq$1,000 nonlegal-state residents within that reach. Capped is the state-year
entry-cap indicator (CO/MT always capped, NM always free, OR switches capped in 2022). State and state
$\times$ quarter fixed effects via a single concatenated state-quarter key (avoids phantom empty cells from
this unbalanced panel). The Capped main effect is exactly collinear with state $\times$ quarter fixed effects
and is not separately estimated (confirmed by matrix-rank check). County-clustered standard errors.
\end{minipage}
\end{table}
